\section{Conclusiones}
\label{sec:conclusiones}

\begin{itemize}
    \item La regresión logística implementada en PyTorch, con evaluación,
    gradiente y entrenamiento propios, separa sin error los datos
    sintéticos y alcanza sobre las características de MPNet una tasa de
    aciertos de $0.9906\pm0.0016$ y un F1 de $0.9641\pm0.0063$ en 30
    particiones. Es el método más sólido evaluado con 30 corridas.
    \item Reducir las 768 características a 384 componentes principales
    conserva el 96.16~\% de la varianza y prácticamente el mismo F1
    ($-0.0004$ en promedio). Con 2 componentes se pierde información útil
    para separar las clases y el F1 baja a 0.7845.
    \item El detector de anomalías por histogramas funciona sin ejemplos de
    spam, pero queda por debajo de la regresión en las 30 particiones (F1
    $0.7043\pm0.0247$ con $D=2$). Su punto débil son los ham poco
    frecuentes que marca como spam. Además, la suma de verosimilitudes de
    muchas componentes favorece al spam: con $D=384$ el detector casi no lo
    encuentra.
    \item Los modelos de lenguaje sin ajuste reconocen casi todo el spam,
    pero producen muchas falsas alarmas; el mayor modelo evaluado, Qwen 14B
    en NF4, obtiene el mejor F1 de los tres (0.7126).
    \item Un ajuste eficiente con LoRA sobre Qwen 14B en NF4 corrige ese
    sesgo: en su prueba reservada, los falsos positivos bajan de 174 a 5 y
    el F1 sube de 0.7120 a 0.9708, también en las filas cuyo texto no
    aparece en el entrenamiento.
    \item La tasa de aciertos sola es engañosa en este dataset desbalanceado:
    predecir siempre ham ya da 0.866. El F1 de la clase spam distingue
    mejor los métodos y refleja el costo de las falsas alarmas.
\end{itemize}

\subsection{Limitaciones}
\label{sec:limitaciones}

\begin{itemize}
    \item Las 30 particiones y las 30 muestras se toman del mismo dataset y
    se solapan; sus desviaciones describen la variación entre particiones,
    no intervalos de confianza de experimentos independientes.
    \item La partición es por filas y conserva los textos repetidos del
    archivo original. En promedio, el 9.7~\% de los mensajes de prueba
    (entre 145 y 181 de 1672) tiene una copia en el entrenamiento, lo que
    puede sobreestimar el desempeño en mensajes nuevos.
    \item El detector usa $D=2$, elegida con validación interna, en lugar
    de reutilizar $D=384$ como pide la lectura literal del apartado 2.1.
    Ambos casos se informan, y la elección de $D$ usó etiquetas de spam de
    validación, por lo que el procedimiento completo no es enteramente no
    supervisado.
    \item Los modelos de lenguaje sin ajuste y el modelo ajustado se evalúan
    con protocolos distintos de los puntos 1 y 2. El ajuste con LoRA es un
    único entrenamiento en una partición, sin estimación de su variabilidad.
    \item No se comprobó si los corpus de preentrenamiento de los modelos de
    lenguaje incluyen esta colección pública, así que no puede descartarse
    exposición previa a los mensajes.
    \item Las comparaciones entre modelos de lenguaje cambian a la vez
    familia, tamaño y precisión de los pesos; no aíslan el efecto de cada
    factor.
\end{itemize}
